% ============================================================================
% REFERENCES
% ============================================================================
\chapter*{References}
\addcontentsline{toc}{chapter}{References}

\begin{enumerate}[label={[\arabic*]}, leftmargin=2cm]

\item Vaswani, A., Shazeer, N., Parmar, N., Uszkoreit, J., Jones, L., Gomez, A.N., Kaiser, L. and Polosukhin, I. (2017). ``Attention Is All You Need.'' \textit{Advances in Neural Information Processing Systems (NeurIPS)}, Vol. 30. This seminal paper introduced the Transformer architecture which eliminated recurrence in favor of pure self-attention mechanisms, achieving state-of-the-art results on WMT English-to-German (28.4 BLEU) and English-to-French (41.0 BLEU) translation benchmarks while enabling massive parallelization during training.

\item Touvron, H., Martin, L., Stone, K., et al. (2023). ``Llama 2: Open Foundation and Fine-Tuned Chat Models.'' \textit{arXiv preprint arXiv:2307.09288}. Meta AI released the Llama 2 family (7B, 13B, 70B parameters) as open-weight models trained on 2T tokens, achieving competitive performance with proprietary models on MMLU (68.9\% for 70B) and introducing RLHF-aligned chat variants that demonstrated safety improvements across 2,000 adversarial prompts.

\item Meta AI (2024). ``Llama 3.3 Model Card.'' Llama-3.3-70B represents the latest iteration of Meta's open model family, trained on 15T+ tokens with GQA (8 KV heads), 128K context window, and RoPE ($\theta = 500{,}000$), achieving 82.0\% on MMLU (5-shot) and 93.0\% on GSM8K, establishing it as the strongest open model for research-grade text generation tasks.

\item Alibaba Cloud (2025). ``Qwen3 Technical Report.'' The Qwen3-32B model employs a Mixture-of-Experts architecture achieving 78.4\% on MMLU with particularly strong performance (85.2\%) on code generation benchmarks (HumanEval+), making it optimal for structured JSON output generation in agentic AI pipelines.

\item Dao, T., Fu, D.Y., Ermon, S., Rudra, A. and R\'{e}, C. (2022). ``FlashAttention: Fast and Memory-Efficient Exact Attention with IO-Awareness.'' \textit{NeurIPS 2022}. This work reduced the memory complexity of self-attention from $O(n^2)$ to $O(n)$ through tiled, IO-aware CUDA kernel implementations, enabling 3$\times$ speedup for BERT-large training and supporting sequences up to 64K tokens on single A100 GPUs.

\item Su, J., Lu, Y., Pan, S., Murtadha, A., Wen, B. and Liu, Y. (2021). ``RoFormer: Enhanced Transformer with Rotary Position Embedding.'' \textit{arXiv:2104.09864}. Introduced Rotary Position Embeddings (RoPE) which encode relative position information through rotation matrices applied to query and key vectors, enabling length generalization beyond training sequence lengths without additional parameters.

\item Wu, Q., Bansal, G., Zhang, J., et al. (2023). ``AutoGen: Enabling Next-Gen LLM Applications via Multi-Agent Conversation.'' \textit{arXiv:2308.08155}, Microsoft Research. Introduced the concept of conversable agents for multi-agent collaboration, demonstrating 23\% improvement in code generation accuracy through agent debate patterns compared to single-agent baselines.

\item Zheng, L., Chiang, W.L., Sheng, Y., et al. (2023). ``Judging LLM-as-a-Judge with MT-Bench and Chatbot Arena.'' \textit{NeurIPS 2023}. Demonstrated that GPT-4's evaluation of text quality achieves 85\% agreement with human expert annotators on the MT-Bench dataset, supporting the feasibility of using LLMs as automated quality assessors in production systems.

\item Karpathy, A. (2024). ``autoresearch --- Autonomous AI Research Agent.'' \textit{GitHub Repository}. Introduced a minimal yet powerful paradigm for autonomous experimentation where an LLM proposes code modifications, executes experiments, evaluates results against baseline metrics, and automatically retains improvements while reverting failures.

\item Goldberg, D.E. and Deb, K. (1991). ``A Comparative Analysis of Selection Schemes Used in Genetic Algorithms.'' \textit{Foundations of Genetic Algorithms}, Vol. 1, pp. 69--93. Formalized tournament selection as a selection operator in evolutionary computation, demonstrating that tournament size $k$ controls selection pressure and that $(1+\lambda)$ strategies provably converge to optimal solutions under mild assumptions.

\item Ainslie, J., Lee-Thorp, J., de Jong, M., Zemlyanskiy, Y., Lebr\'{o}n, F. and Sanghai, S. (2023). ``GQA: Training Generalized Multi-Query Transformer Models from Multi-Head Checkpoints.'' \textit{EMNLP 2023}. Grouped-Query Attention reduces KV-cache memory by 4--8$\times$ while retaining 99.5\% of Multi-Head Attention quality, enabling efficient inference for large-scale language models.

\item Wooldridge, M. and Jennings, N.R. (1995). ``Intelligent Agents: Theory and Practice.'' \textit{The Knowledge Engineering Review}, 10(2), pp. 115--152. Foundational work defining the BDI (Belief-Desire-Intention) model for autonomous agents, establishing the theoretical framework for multi-agent systems that influenced subsequent developments including modern LLM-based agent architectures.

\item Gu, A. and Dao, T. (2023). ``Mamba: Linear-Time Sequence Modeling with Selective State Spaces.'' \textit{arXiv:2312.00752}. Introduced selective state space models achieving Transformer-quality language modeling at $O(n)$ complexity instead of $O(n^2)$, processing 5$\times$ faster than FlashAttention-2 on A100 GPUs for sequences beyond 8K tokens.

\item Peng, B., Alcaide, E., Anthony, Q., et al. (2023). ``RWKV: Reinventing RNNs for the Transformer Era.'' \textit{EMNLP 2023 Findings}. Proposed a novel architecture combining Transformer-level parallelism during training with RNN-level efficiency during inference, achieving competitive perplexity (16.2 on WikiText-103) with constant memory requirements independent of sequence length.

\item Hunter, J.D. (2007). ``Matplotlib: A 2D Graphics Environment.'' \textit{Computing in Science \& Engineering}, 9(3), pp. 90--95. The foundational Python visualization library that established the paradigm for programmatic scientific chart generation, serving as the basis upon which higher-level visualization frameworks (Seaborn, Plotly) were subsequently developed.

\end{enumerate}
